% =====================================================================
%  ACADEMIC REPORT  --  template for inscrive.io
%  ---------------------------------------------------------------
%  Engine   : pdfLaTeX          Bibliography : biber
%  Structure: main.tex holds metadata only. Text lives in content/,
%             all design lives in style/universityreport.sty.
%
%  To make this yours:
%    1. Replace figures/logo-placeholder.pdf with your institution logo.
%    2. Edit the brand colours in style/universityreport.sty (section 1).
%    3. Fill in the metadata block below and start writing in content/.
% =====================================================================
\documentclass[12pt,a4paper,oneside]{report}

\usepackage{style/universityreport}
\addbibresource{references.bib}

% ---------------------------------------------------------------------
%  METADATA  --  everything on the cover page comes from here
% ---------------------------------------------------------------------
\reporttitle      {Measuring Compilation Latency in Collaborative Document Systems}
\reportsubtitle   {A reproducible benchmark of build pipelines for large multi-author projects}
\reportauthor     {A. N. Author}
\reportinstitution{Your University \textperiodcentered\ Faculty or Department}
\reportmeta       {Supervisor: Prof.\ B. C. Supervisor}
\reportdate       {\today}
\reportlogo       {figures/logo-placeholder.pdf}

\begin{document}

\makecoverpage

% ---------------------------------------------------------------------
\frontmatterpages
% ---------------------------------------------------------------------
%  Abstract  --  unnumbered, but listed in the table of contents
% ---------------------------------------------------------------------
\chapter*{Abstract}
\addcontentsline{toc}{chapter}{Abstract}
\markboth{Abstract}{}

\begingroup
\setlength{\leftskip}{0pt}\noindent
Collaborative writing platforms rebuild a document every time an author
pauses typing, which makes compilation latency the single most visible
performance characteristic of the system. This report presents a
reproducible benchmark of three build strategies -- a cold baseline, a
warm package cache, and incremental recompilation -- measured across
documents of \numrange{20}{160} pages. We find that incremental
recompilation keeps median build time below \qty{1.4}{\second} even for
the largest documents, an order of magnitude faster than the cold
baseline, and that the crossover point where caching stops paying off
occurs at roughly 60 pages. We close with a short discussion of what
these numbers mean for multi-author projects and how the measurement
harness can be reused.
\endgroup

\vspace{6mm}
\noindent\textbf{\sffamily Keywords:}\enspace
compilation latency \textperiodcentered\ reproducible builds
\textperiodcentered\ collaborative editing \textperiodcentered\
document engineering

\tableofcontents
\cleardoublepage

% ---------------------------------------------------------------------
\mainmatterpages
\chapter{Introduction}
\label{chap:introduction}

Typesetting systems are usually evaluated on the quality of their
output. For a collaborative editor, however, the more decisive property
is how long an author waits between pressing a key and seeing the
result. That wait is what determines whether a writing session feels
like editing a document or like submitting a job to a queue.

\section{Motivation}

A modern research document is rarely a single file written by a single
person. It is a tree of included sections, a bibliography that is
synchronised from a reference manager, a directory of generated figures,
and a version history shared between several authors
\autocite{wilson2014best}. Every one of those moving parts adds work to
the build, and the build runs far more often than anyone reads the
finished PDF.

\begin{keybox}[Why latency dominates]
In a two-hour writing session an author triggers on the order of
\num{200} rebuilds but reads the full output perhaps twice. Reducing
median build time from \qty{8}{\second} to \qty{1}{\second} gives back
roughly \qty{23}{\minute} of that session.
\end{keybox}

\section{Scope and contributions}

This report is deliberately narrow. It measures build time and nothing
else: no rendering fidelity, no editor responsiveness, no network
effects. Within that scope it makes three contributions.

\begin{enumerate}
  \item A benchmark harness that produces documents of controlled size
        and structure, described in \cref{sec:harness}.
  \item A comparison of three build strategies over eight document
        sizes, reported in \cref{chap:results}.
  \item A cost model, \cref{eq:model}, that predicts build time from
        page count and cache state within \qty{9}{\percent}.
\end{enumerate}

\section{How to read this report}

\Cref{chap:background} fixes notation and states the assumptions behind
the cost model. \Cref{chap:method} describes the harness and the
measurement protocol. \Cref{chap:results} presents the measurements, and
\cref{chap:conclusion} discusses what generalises beyond our setup.
Supporting material is in \cref{chap:appendix-a}.

\begin{remark}
Every number in this report was produced by the scripts in the
repository that accompanies it. Re-running \texttt{make bench} should
reproduce \cref{fig:latency} up to hardware differences.
\end{remark}

\section{Conventions}

Symbols are introduced where they first appear and collected in
\cref{tab:notation}. Times are wall-clock seconds measured inside the
build worker; all other quantities are dimensionless.

\begin{table}[htbp]
  \centering
  \caption{Notation used throughout the report.}
  \label{tab:notation}
  \small
  \begin{tabularx}{\textwidth}{@{}c L l@{}}
    \toprule
    \textbf{\sffamily Symbol} & \textbf{\sffamily Meaning} & \textbf{\sffamily Unit} \\
    \midrule
    $p$      & page count of the compiled document        & pages \\
    $n$      & number of source files in the project      & files \\
    $K$      & number of typesetting passes per build     & --- \\
    $\alpha$ & fixed initialisation cost of a build       & s \\
    $\beta$  & marginal typesetting cost per page         & s/page \\
    $\gamma$ & per-pass overhead independent of $p$       & s \\
    $\rho$   & fraction of pages changed since the last build & --- \\
    \bottomrule
  \end{tabularx}
\end{table}

A short note on the markup this template expects:

\begin{description}
  \item[Cross-references] Written with \texttt{cleveref}. Typing
        \verb|\cref{tab:notation}| gives \cref{tab:notation}, with the
        word \enquote{Table} supplied automatically, so renaming a
        float never leaves a stale label behind.
  \item[Citations] Written with \texttt{biblatex}. Typing
        \verb|\autocite{key}| produces \autocite{perez2007ipython}; use
        \verb|\textcite{key}| instead when the author should be part of
        the sentence, as in \textcite{hunter2007matplotlib}.
  \item[Units and numbers] Written with \texttt{siunitx}, so
        \verb|\qty{1.4}{\second}| produces \qty{1.4}{\second} and stays
        consistent across the whole document.
\end{description}

\chapter{Background and cost model}
\label{chap:background}

\section{Notation}

Let $D$ be a document composed of $n$ source files and let $p$ denote
its page count. A build is a function that maps $D$ to a PDF; we write
$T(D)$ for its wall-clock duration. The build is decomposed into a fixed
cost $\alpha$ (format loading, package initialisation), a per-page
typesetting cost $\beta$, and a per-run overhead $\gamma_k$ for the
$k$-th pass.

\begin{definition}[Warm cache]
\label{def:warm}
A build is \emph{warm} if every package, font and image required by $D$
is already materialised on the worker that executes it. Otherwise the
build is \emph{cold}.
\end{definition}

\section{A two-term model}

Under \cref{def:warm} the total time of a $K$-pass build is the sum of
the per-pass costs,
%
\begin{equation}
  \label{eq:model}
  T(p, K) \;=\; \alpha \;+\; \sum_{k=1}^{K}\bigl(\gamma_k + \beta_k\, p\bigr)
  \;\approx\; \alpha + K\gamma + \beta p ,
\end{equation}
%
where the approximation assumes that the per-page cost is stable across
passes. The assumption holds well in practice because the second and
later passes differ from the first only in the resolution of forward
references.

For incremental builds the model changes shape. Only the fraction
$\rho \in [0,1]$ of pages touched since the previous build is
re-typeset, so
%
\begin{align}
  T_{\mathrm{inc}}(p, K) &= \alpha_{\mathrm{inc}} + K\gamma + \beta \rho p ,
  \label{eq:inc} \\[2pt]
  \rho &= \frac{\lvert \mathcal{P}_{\text{dirty}} \rvert}{p} ,
  \qquad
  \mathcal{P}_{\text{dirty}} = \bigl\{\, \pi \in D : h(\pi) \neq h'(\pi) \,\bigr\},
  \label{eq:rho}
\end{align}
%
with $h$ and $h'$ the content hashes before and after the edit.

\begin{theorem}[Crossover]
\label{thm:crossover}
Caching outperforms incremental recompilation only while
$\alpha - \alpha_{\mathrm{inc}} > \beta p\,(1-\rho)$. For fixed $\rho$
the right-hand side grows linearly in $p$, so there exists a page count
$p^{\star}$ beyond which incremental recompilation is always faster.
\end{theorem}

\begin{proof}
Subtract \cref{eq:inc} from \cref{eq:model} and solve the resulting
inequality for $p$. The $K\gamma$ terms cancel.
\end{proof}

\begin{lemma}
\label{lem:bound}
If $\rho \le \tfrac{1}{2}$ then $T_{\mathrm{inc}} \le \tfrac{1}{2}T + \alpha_{\mathrm{inc}}$.
\end{lemma}

\begin{keybox}[What the model predicts]
\Cref{thm:crossover} says the advantage of a warm cache is a constant,
while the advantage of incremental work grows with document size. Long
documents should therefore favour incremental builds even when the cache
is perfectly warm. \Cref{chap:results} tests exactly this.
\end{keybox}

\section{Related work}

Build-time modelling for document systems borrows directly from the
compiler literature, and the canonical description of the multi-pass
structure that makes \cref{eq:model} necessary remains
\autocite{knuth1984texbook,lamport1994latex}. Practical package-level
detail is collected in \autocite{mittelbach2023companion}. Scheduling of
shared build workers, which we treat as out of scope, is addressed by
\textcite{example2026conference}.

\chapter{Method}
\label{chap:method}

\section{The benchmark harness}
\label{sec:harness}

The harness generates synthetic documents with a controlled mix of
prose, mathematics, tables and vector figures, compiles them under each
build strategy, and records wall-clock time. \Cref{fig:pipeline} shows
the stages.

\begin{widefigure}
  \centering
  \begin{tikzpicture}[
      node distance=10mm and 11mm,
      box/.style={draw=hairline,fill=white,rounded corners=2pt,
                  minimum height=10mm,minimum width=25mm,align=center,
                  font=\sffamily\scriptsize,line width=0.6pt,inner sep=4pt},
      hot/.style={box,draw=accent,line width=1pt,text=ink},
      lbl/.style={font=\sffamily\scriptsize\itshape,text=muted,align=left},
      arr/.style={-{Stealth[length=4pt]},draw=muted,line width=0.6pt},
      feedback/.style={arr,dashed,dash pattern=on 2pt off 1.6pt}
    ]
    \node[box]               (gen)   {Document\\generator};
    \node[box,right=of gen]  (src)   {Source tree\\$n$ files};
    \node[hot,right=of src]  (build) {Build worker\\$K$ passes};
    \node[box,right=of build](pdf)   {PDF +\\timing log};

    \node[box,below=of src]   (hash)  {Content\\hashes};
    \node[box,below=of build] (cache) {Package\\cache};

    \begin{scope}[on background layer]
      \node[fill=surface,rounded corners=3pt,fit=(hash)(cache),inner sep=7pt] (panel) {};
    \end{scope}

    \draw[arr] (gen)   -- (src);
    \draw[arr] (src)   -- (build);
    \draw[arr] (build) -- (pdf);
    \draw[arr] (hash)  -- (src);
    \draw[arr] (cache) -- (build);

    % the worker writes fresh hashes back after every run
    \coordinate (turn) at ($(panel.south)+(0,-7mm)$);
    \draw[feedback] (pdf.south) |- (turn) -| (hash.south);

    \node[lbl,anchor=north] at ($(turn)+(0,-1mm)$) {new hashes written back after every run};
    \node[lbl,anchor=east,align=right] at ($(panel.west)+(-3mm,0)$) {reused\\between runs};
  \end{tikzpicture}
  \caption[Benchmark pipeline]{Stages of the benchmark harness. Only the
           build worker is timed; generation and hashing happen outside
           the measured window.}
  \label{fig:pipeline}
\end{widefigure}

\section{Build strategies}

Three strategies were compared. They differ only in what the worker is
allowed to reuse between runs.

\begin{table}[htbp]
  \centering
  \caption{The three build strategies under test.}
  \label{tab:strategies}
  \small
  \begin{tabularx}{\textwidth}{@{}l L l@{}}
    \toprule
    \textbf{\sffamily Strategy} & \textbf{\sffamily Reuses} & \textbf{\sffamily Passes $K$} \\
    \midrule
    Baseline     & nothing; a fresh container per run            & 3 \\
    Cached       & package tree, fonts, generated images          & 3 \\
    Incremental  & all of the above plus page-level content hashes & 1--3 \\
    \bottomrule
  \end{tabularx}
\end{table}

\section{Measurement protocol}

Each configuration was run \num{30} times. The first three runs are
discarded as warm-up, and we report the median of the remaining
\num{27}. Timing uses a monotonic clock taken inside the worker, so
scheduling delay in the queue is excluded.

\begin{lstlisting}[language=Python,caption={Core of the timing loop.},label={lst:timing}]
def measure(doc, strategy, repeats=30, warmup=3):
    """Return the median wall-clock build time in seconds."""
    samples = []
    for i in range(repeats):
        worker = Worker(strategy)          # fresh state for "baseline"
        t0 = time.monotonic()
        worker.build(doc)
        samples.append(time.monotonic() - t0)
    return statistics.median(samples[warmup:])
\end{lstlisting}

\begin{remark}
Using the median rather than the mean matters here: the distribution has
a long right tail caused by occasional font-cache regeneration, and a
single outlier can move the mean by more than \qty{20}{\percent}.
\end{remark}

\chapter{Results}
\label{chap:results}

\section{Build time against document size}

\sidenote{The raw timings live in \texttt{measurements.csv} under
\texttt{data/}. Edit that file and every figure in this chapter
follows.}%
\Cref{fig:latency} plots median build time against page count for the
three strategies. The data is read directly from
\texttt{data/measurements.csv}, so the figure updates whenever the
benchmark is re-run.

\begin{widefigure}
  \centering
  \begin{tikzpicture}
    \begin{axis}[
        width=0.94\fullwidth, height=62mm,
        xlabel={Document size (pages)},
        ylabel={Median build time (s)},
        xmin=15, xmax=165, ymin=0,
        grid=major,
        grid style={hairline,line width=0.4pt},
        axis line style={muted,line width=0.5pt},
        tick label style={font=\sffamily\scriptsize,color=muted},
        label style={font=\sffamily\small,color=ink},
        legend style={font=\sffamily\scriptsize,draw=hairline,
                      fill=white,at={(0.02,0.97)},anchor=north west,
                      cells={anchor=west}},
        every axis plot/.append style={line width=1.1pt,mark size=1.6pt},
      ]
      \addplot[accent,mark=*]        table[x=pages,y=baseline,   col sep=comma]{data/measurements.csv};
      \addplot[accentdark,mark=square*,dashed] table[x=pages,y=cached,col sep=comma]{data/measurements.csv};
      \addplot[muted,mark=triangle*] table[x=pages,y=incremental,col sep=comma]{data/measurements.csv};
      \legend{Baseline, Cached, Incremental}
    \end{axis}
  \end{tikzpicture}
  \caption[Median build time by strategy]{Median build time against
    document size. The baseline grows super-linearly once the document
    exceeds the worker's memory budget at roughly \num{120} pages;
    incremental builds stay almost flat.}
  \label{fig:latency}
\end{widefigure}

\section{Where the time goes}

Splitting a single \num{100}-page build into its phases
(\cref{fig:breakdown}) shows that the fixed cost $\alpha$ of
\cref{eq:model} is dominated by package initialisation, which is exactly
the part a warm cache removes.

\begin{widefigure}
  \centering
  \begin{subfigure}[t]{0.48\fullwidth}
    \centering
    \begin{tikzpicture}
      \begin{axis}[
          width=\textwidth, height=48mm,
          ybar, bar width=9pt,
          symbolic x coords={Init,Pass 1,Pass 2,Pass 3},
          xtick=data, ymin=0, ymax=4.2,
          ylabel={Time (s)},
          nodes near coords, nodes near coords style={font=\sffamily\tiny,color=muted},
          grid=major, grid style={hairline,line width=0.4pt},
          axis line style={muted,line width=0.5pt},
          tick label style={font=\sffamily\scriptsize,color=muted},
          label style={font=\sffamily\scriptsize,color=ink},
        ]
        \addplot[draw=accent,fill=accent!22] coordinates
          {(Init,3.61) (Pass 1,2.14) (Pass 2,1.18) (Pass 3,1.01)};
      \end{axis}
    \end{tikzpicture}
    \caption{Cold baseline.}
    \label{fig:breakdown-cold}
  \end{subfigure}
  \hfill
  \begin{subfigure}[t]{0.48\fullwidth}
    \centering
    \begin{tikzpicture}
      \begin{axis}[
          width=\textwidth, height=48mm,
          ybar, bar width=9pt,
          symbolic x coords={Init,Pass 1,Pass 2,Pass 3},
          xtick=data, ymin=0, ymax=4.2,
          ylabel={Time (s)},
          nodes near coords, nodes near coords style={font=\sffamily\tiny,color=muted},
          grid=major, grid style={hairline,line width=0.4pt},
          axis line style={muted,line width=0.5pt},
          tick label style={font=\sffamily\scriptsize,color=muted},
          label style={font=\sffamily\scriptsize,color=ink},
        ]
        \addplot[draw=accentdark,fill=accentdark!22] coordinates
          {(Init,0.34) (Pass 1,1.96) (Pass 2,0.68) (Pass 3,0.54)};
      \end{axis}
    \end{tikzpicture}
    \caption{Warm cache.}
    \label{fig:breakdown-warm}
  \end{subfigure}
  \caption[Phase breakdown]{Phase breakdown of a \num{100}-page build.
    The cache removes almost all of the initialisation cost but leaves
    the typesetting passes untouched.}
  \label{fig:breakdown}
\end{widefigure}

\section{Fitted parameters}

Fitting \cref{eq:model} to the measured series by ordinary least squares
gives the parameters in \cref{tab:fit}. The residual standard error is
below \qty{9}{\percent} for all three strategies, which is within the
run-to-run spread of the measurements themselves.

Two things stand out. First, the fixed cost $\alpha$ falls by more than
an order of magnitude as soon as the package tree is cached, confirming
the phase breakdown in \cref{fig:breakdown}. Second, the per-page cost
$\beta$ of the incremental strategy is not zero: even an unchanged page
must be re-shipped to the output driver, and that residual cost is what
keeps \cref{lem:bound} from being tight.

\begin{table}[htbp]
  \centering
  \caption{Least-squares fit of \cref{eq:model} to the measured series.}
  \label{tab:fit}
  \small
  \begin{tabular}{@{}l S[table-format=1.3] S[table-format=1.4] S[table-format=1.3] S[table-format=2.1]@{}}
    \toprule
    \textbf{\sffamily Strategy} & {$\alpha$ (s)} & {$\beta$ (s/page)} & {$\gamma$ (s)} & {RSE (\%)} \\
    \midrule
    Baseline    & 3.610 & 0.0742 & 0.284 & 8.4 \\
    Cached      & 0.340 & 0.0301 & 0.271 & 6.1 \\
    Incremental & 0.318 & 0.0061 & 0.119 & 4.7 \\
    \bottomrule
  \end{tabular}
\end{table}

\begin{keybox}[Crossover confirmed]
The measured crossover between the cached and incremental strategies
falls at $p^{\star} \approx 58$ pages, within \qty{4}{\percent} of the
value that \cref{thm:crossover} predicts from the fitted parameters in
\cref{tab:fit}.
\end{keybox}

\chapter{Discussion and conclusion}
\label{chap:conclusion}

\section{What generalises}

The absolute numbers in \cref{chap:results} belong to one hardware
configuration and should not be quoted outside it. The \emph{shape} of
the curves is more robust: a fixed initialisation cost that caching
removes, a per-page cost that only incremental work removes, and a
crossover between the two that moves with document size exactly as
\cref{thm:crossover} predicts.

\section{Threats to validity}

\begin{description}
  \item[Synthetic documents.] Generated sources have a more uniform mix
        of content than real manuscripts. Documents dominated by large
        raster figures behave differently, because image handling is not
        page-proportional.
  \item[Single worker.] All measurements were taken on an unloaded
        worker. Contention between concurrent builds is left to future
        work \autocite{example2026conference}.
  \item[Median reporting.] Medians hide the tail that users actually
        complain about. \Cref{tab:percentiles} in \cref{chap:appendix-a}
        gives the 95th percentile for completeness.
\end{description}

\section{Conclusion}

Incremental recompilation is the only strategy of the three that keeps
build time roughly independent of document size, and it becomes the
better choice at around \num{60} pages -- well below the size of a
typical thesis. For collaborative systems, where the build runs hundreds
of times per session, that is the difference between a tool that keeps
up with the author and one that does not.

\begin{proposition}
Under the fitted parameters of \cref{tab:fit} and $\rho \le 0.1$, a
\num{300}-page document rebuilds incrementally in under \qty{2.5}{\second}.
\end{proposition}


% ---------------------------------------------------------------------
\printbibliography[heading=bibintoc,title={References}]

% ---------------------------------------------------------------------
\appendix
\chapter{Supplementary material}
\label{chap:appendix-a}

\section{Tail latency}

\begin{table}[htbp]
  \centering
  \caption{Build time percentiles for a \num{100}-page document.}
  \label{tab:percentiles}
  \small
  \begin{tabular}{@{}l S[table-format=2.2] S[table-format=2.2] S[table-format=2.2]@{}}
    \toprule
    \textbf{\sffamily Strategy} & {p50 (s)} & {p90 (s)} & {p95 (s)} \\
    \midrule
    Baseline    & 7.94 & 11.62 & 14.08 \\
    Cached      & 3.52 &  4.91 &  6.03 \\
    Incremental & 0.94 &  1.28 &  1.51 \\
    \bottomrule
  \end{tabular}
\end{table}

\section{Reproducing the measurements}

\begin{lstlisting}[language=bash,caption={Full reproduction sequence.},label={lst:repro}]
git clone https://github.com/<your-org>/<this-repo>.git
cd <this-repo>
make bench          # writes data/measurements.csv
latexmk -pdf main   # rebuilds this report from the new data
\end{lstlisting}

\begin{remark}
The figures in \cref{chap:results} read \texttt{data/measurements.csv}
at compile time, so no figure needs to be regenerated by hand after a
new benchmark run.
\end{remark}


\end{document}
